\section{Selection of the entrance law}\label{sec:entrance-selection}

Proposition~\ref{prop:gaussian-subsequence} gives positive-time density
families for the edge semigroup, but an extracted family might still depend
on randomness beyond the edge points.  The uniform-start estimate of
Proposition~\ref{prop:uniform-seed-escape} is averaged over the represented
environment.  We first prove fixed-\(g\) uniqueness from a sequential
escape criterion and derive the zero-coordinate boundary from that
criterion.  We then extract a qualifying sequence for almost every
represented environment and obtain a family measurable in \(g\).

Fix \(g\) in a common Borel full-measure subset of the edge set chosen
in Proposition~\ref{prop:uniform-seed-escape} on which the conclusions
about the conditioned law and the closed gradient form in Items 1 and 3
of Proposition~\ref{prop:stationary-inputs} hold with their stated
source versions.  The rate, finiteness, and positivity conclusions of
Proposition~\ref{prop:uniform-seed-escape} hold on this set as well.
In this section write
\[
 r_i=g_i+b(g),\qquad D=D_g(b(g)),\qquad \Pi=\Pi_g,\qquad H=H_g.
\]
Thus \(D>0\), and there are constants \(0<c_g<C_g<\infty\) such that
\begin{equation}\label{entr:rates}
 c_g(1+i)^{2/3}\le r_i\le C_g(1+i)^{2/3}\qquad(i\ge1).
\end{equation}
Constants with a subscript \(g\) below may be enlarged or decreased
and depend only on this fixed realization, unless other dependencies
are displayed.
Recall the finite, strictly positive random variable
\begin{equation}\label{entr:mass}
 M(x)=\sum_{i\ge1}\frac{x_i^2}{2r_i},\qquad x\sim\Pi.
\end{equation}
Its finiteness follows from the uniform coordinate moment bound
\(\Pi[x_i^2]\le C_g/r_i\) and \(\sum_i r_i^{-2}<\infty\); its positivity
follows from the absolute continuity of a one-coordinate marginal.
These properties of \(\Pi\) are part of the conditioned edge law
\cite[Proposition~2.3]{FUA}.

\begin{theorem}[Entrance selection]\label{thm:entrance-selection}
For every \(g\) as above, there is at most one family
\((h_s)_{s>0}\) of probability densities with respect to \(\Pi_g\) with
the following properties:
\begin{enumerate}
 \item \(h_s\ge0\), \(\Pi_g[h_s]=1\), and
 \(\sup_{s_0\le s\le s_1}\|h_s\|_{L^\infty(\Pi_g)}<\infty\) whenever
 \(0<s_0<s_1<\infty\);
 \item \(h_{s+t}=S_t^g h_s\) in \(L^2(\Pi_g)\) for \(s,t>0\), where
 \(S_t^g=\exp(-t c_*H_g)\);
 \item there is a sequence \(s_k\downarrow0\) such that, for every
 \(i\ge1\) and \(\varepsilon>0\),
 \begin{equation}\label{entr:escape}
 \Pi_g\left[h_{s_k}\,
       \ind_{\{x_i^2/M(x)>\varepsilon\}}\right]\longrightarrow0.
 \end{equation}
\end{enumerate}
Any such family satisfies
\begin{equation}\label{entr:zero-boundary}
 h_s\Pi_g\ \weakto\ \delta_{\boldsymbol 0}
 \quad\text{on \(\R^{\N}\) with its product topology, as \(s\downarrow0\)}.
\end{equation}
For almost every \(g\) such a family exists.
It can be chosen measurably in \(g\), as a family of \(L^2(\Pi_g)\)
density classes, and with jointly measurable density representatives in
\((g,s,x)\).  Every represented Gaussian limit in
Proposition~\ref{prop:gaussian-subsequence} is this family almost surely.
\end{theorem}

The hypothesis \eqref{entr:escape} describes information retained from
the finite uniform initialization.  Different candidate families may
use different sequences in Item 3.  The zero-coordinate limit
\eqref{entr:zero-boundary} is a consequence of that hypothesis and
holds as \(s\) decreases to zero through all positive times.

We prove uniqueness in the time variable
\begin{equation}\label{entr:unit-time}
 u=c_*s,\qquad \widetilde h_u=h_{u/c_*},\qquad
 \widetilde S_u=e^{-uH}.
\end{equation}
Then \(\widetilde h_{u+v}=\widetilde S_v\widetilde h_u\).
The sequence \(u_k=c_*s_k\) has exactly the escape property
\eqref{entr:escape}.  All elapsed times in the next three subsections
refer to \(\widetilde S\).  We return to \(s=u/c_*\) at the end.

\subsection{Angular coordinates and a common diffusion factor}

We need a stochastic realization only after a positive time, where the
initial law has a bounded density.  Put \(m=12\) and define
\[
 R(x)=\left(\sum_{j\le m}x_j^2\right)^{1/2}.
\]
On \(\{R>0\}\), set
\[
 \theta_i(x)=\frac{x_i}{R(x)},\qquad
 \overline r(x)=\sum_{j\le m}r_j\theta_j(x)^2 .
\]
The functions \(\theta_i\) are homogeneous of degree zero: if a finite
set \(J\) contains the head \(\{1,\ldots,m\}\) and \(i\), then
\(\sum_{j\in J}x_j\partial_j\theta_i=0\).
In the finite-coordinate density of \(\Pi\), the factor imposing the
square constraint has a radial derivative.  Its contribution to the
generator therefore vanishes on \(\theta_i\).  The next lemma justifies
this cancellation for the associated diffusion, including avoidance
of \(R=0\).

\begin{lemma}[Angular realization]\label{entr:angular-realization}
Let \(w\) be a bounded probability density with respect to \(\Pi\).
There is a continuous process \(X(t)=(X_i(t))_{i\ge1}\) on
\(\R^\N\), with law \((\widetilde S_t w)\Pi\) at time \(t\), and
independent standard Brownian motions \((W_i)_{i\ge1}\) in its filtration,
with the following property.  Almost surely \(R(X(t))>0\) for all
\(t\ge0\).  With \(R(t)=R(X(t))\),
\(\Theta_i(t)=\theta_i(X(t))\), and
\(\overline r(t)=\overline r(X(t))\), one has
\begin{align}
 d\Theta_i(t)
 &=
 \frac{\sqrt2}{R(t)}
 \left(dW_i(t)-\Theta_i(t)\sum_{j\le m}\Theta_j(t)\,dW_j(t)\right)
 \notag\\
 &\quad+
 \Theta_i(t)\left(
  \frac{3-m-2\ind_{\{i\le m\}}}{R(t)^2}-r_i+\overline r(t)
 \right)dt .                                      \label{entr:angular-sde}
\end{align}
All identities hold simultaneously for \(i\ge1\).  The Brownian
increments after zero are independent of the initial sigma field.
\end{lemma}

\begin{proof}
The closed form of \(H\) is
\(\cE_g(F,G)=\Pi[\nabla F\cdot\nabla G]\).  Its domain is the closure
of the smooth compact cylinder functions, contains the constants and all
coordinates, and satisfies \(\nabla x_i=e_i\)
\cite[Proposition~5.1]{FUA}.
Smooth contractions on the core extend by closure, so this is a
Dirichlet form.  The same approximation, with smooth clipping to preserve
a common bound, gives the chain and bounded product rules.  In particular
the gradient vanishes almost everywhere on a set on which the function is
constant: apply the chain rule to smooth truncations in successively
smaller neighborhoods of the constant, and use closedness.  These rules
show strong locality.

Here are the topological hypotheses for the diffusion theorem.  The
state space \(\R^\N\) is Polish in the product topology.  The smooth
compact cylinder core consists of continuous functions and contains
a countable family separating its points.  To obtain a compact nest,
for each coordinate choose a smooth
function \(v_j\), equal to one outside a sufficiently large coordinate
interval and zero on a smaller interval, with
\(\|v_j\|_{\cE_g,1}\le\epsilon2^{-j}\).  Such a choice is possible:
coordinate tails have vanishing \(\Pi\)-mass, while a transition of
increasing width makes the gradient norm vanish.  The sum
\(\sum_jv_j\) converges in the form norm and is at least one on the
complement of a compact product box.  The definition of \(1\)-capacity
therefore bounds that complement's capacity by \(\epsilon^2\).
Taking \(\epsilon\downarrow0\) and enlarging successive boxes gives a
compact nest.  The form is thus quasi-regular.  Since \(1\) has zero
energy, its semigroup fixes \(1\); the associated process has infinite
lifetime under \(\Pi\), and hence under \(w\Pi\).

The diffusion theorem and locality criterion for quasi-regular
symmetric Dirichlet forms now give a continuous associated process
\cite[Theorems~IV.3.5 and V.1.11]{MaRockner}.  Its Fukushima
decomposition is supplied by \cite[Theorem~VI.2.5]{MaRockner}.
In particular, the coordinate martingales \(M^i\) are continuous and
have brackets
\[
 \langle M^i,M^j\rangle_t
   =2\int_0^t \nabla x_i\cdot\nabla x_j(X(v))\,dv
   =2\delta_{ij}t.
\]
The bracket identity is the energy-measure identity for the form,
given by \cite[Theorem~5.2.3 and Equation~(5.2.34)]{FOT} after the
quasi-regular regularization of \cite[Chapter~VI]{MaRockner}.
To apply its bounded formula, first take a bounded smooth clip
\(u_N=c_N(x_i)\).  For every bounded form-domain test \(f\), the chain
and product rules give
\[
 2\cE_g(u_Nf,u_N)-\cE_g(u_N^2,f)
       =2\Pi[f|c_N'(x_i)|^2].
\]
Choose the clips with \(|c_N'|\le1\) and \(c_N(x_i)\to x_i\) in form
norm.  Passing to the limit in the energy measures and polarizing gives
the displayed coordinate brackets.
The multidimensional Lévy characterization, for every finite set of
indices, gives \(M^i=\sqrt2W_i\) with the asserted countable family of
Brownian motions.  Brownian motion here is relative to the filtration of
the process, so its future increments are independent of the initial
sigma field.  The usual form exceptional set has \(\Pi\)-measure zero;
it can be discarded for \(w\Pi\).  These facts also give continuous
coordinate paths.

We verify the generator used for \(\theta_i\).  The finite-coordinate
density formula for \(\Pi\), on any finite set \(J\), is
\begin{equation}\label{entr:cylinder-density}
 \frac{q_{J^c}\!\left(D-\sum_{j\in J}(x_j^2-r_j^{-1})\right)}
      {q(D)}
 \prod_{j\in J}\left(\frac{r_j}{2\pi}\right)^{1/2}
                     e^{-r_jx_j^2/2}.
\end{equation}
Here \(q\) is the density of the full centered Gaussian square sum,
and \(q_{J^c}\) is the density after the coordinates in \(J\) are
deleted
\cite[Proposition~2.3]{FUA}.
Both densities exist, \(q(D)>0\), and \(q_{J^c}\) has bounded
derivatives of any prescribed finite order.  For the last assertion,
leave sufficiently many fixed indices outside \(J\); the product of
their characteristic-function moduli is bounded by
\(\prod (1+4t^2r_j^{-2})^{-1/4}\), which is integrable even after
multiplication by the prescribed power of \(|t|\).  Fourier inversion
then gives the assertion.  In particular \eqref{entr:cylinder-density}
is bounded by a constant times the product Gaussian density.

The set \(R=0\) is polar.  Indeed a smooth cutoff equal to one on
\(\{R\le\epsilon\}\), zero on \(\{R\ge2\epsilon\}\), and with derivative
\(O(\epsilon^{-1})\), has squared form norm
\(O(\epsilon^m+\epsilon^{m-2})\), by \eqref{entr:cylinder-density}.
Letting \(\epsilon\downarrow0\) proves zero capacity.  A diffusion
starting outside its exceptional set never hits a polar set, which
proves the stated pathwise positivity of \(R\).

For a cylinder test whose index set \(J\) contains
\(\{1,\ldots,m,i\}\), integrate by parts in
\eqref{entr:cylinder-density}.  The derivative of the factor \(q_{J^c}\)
is radial in the \(J\) coordinates.  Its pairing with the gradient of
\(\theta_i\) vanishes, because
\(\sum_{j\in J}x_j\partial_j\theta_i=0\).  Consequently the
nonpositive generator on this function is
\[
 \cL\theta_i=\Delta\theta_i-\sum_{j\in J}r_jx_j\partial_j\theta_i.
\]
This integration by parts is legitimate across \(R=0\).  More
explicitly, cut off at \(R=\epsilon\) and at large finite-coordinate
radius.  The function and its first two derivatives near zero are
bounded by
\(C_i(1+|x_i|)(1+R^{-3})\).  The products with derivatives of the
scale-\(\epsilon\) cutoff satisfy the same bound on \(R\asymp\epsilon\).
Their squared integrals near zero vanish for \(m=12>6\).
Terms with a derivative on \(q_{J^c}\) are bounded by the same Gaussian
integrals with polynomial factors, since its first derivative is
bounded.  Gaussian tails handle the outer cutoff.  Thus the cutoffs
converge in the form norm, proving \(\theta_i\) belongs to the form
domain, and the weak generator pairing passes to the limit.
The displayed \(\cL\theta_i\) is in \(L^2(\Pi)\) by the same bounds.
Extending the pairing from the core shows
\(\theta_i\in\operatorname{Dom}(H)\) and \(H\theta_i=-\cL\theta_i\).

A direct differentiation gives, with the head vector extended by zeros,
\[
 \partial_j\theta_i=\frac{\delta_{ij}-\theta_i\theta_j\ind_{\{j\le m\}}}{R},
 \qquad
 \cL\theta_i
 =\theta_i\left(\frac{3-m-2\ind_{\{i\le m\}}}{R^2}
                         -r_i+\overline r\right).
\]
The Fukushima drift is the time integral of this \(L^2\) generator.
Its martingale has brackets with the coordinate martingales equal to
\(2\int_0^t\partial_j\theta_i(X(v))\,dv\), and self-bracket
\(2\int_0^t|\nabla\theta_i(X(v))|^2\,dv\).
It therefore equals
\(\sqrt2\sum_j\int_0^t\partial_j\theta_i(X(v))\,dW_j(v)\):
the difference has zero bracket.  Only the head and index \(i\)
occur in this finite sum.  Integrability follows, for instance, from
\[
 \E\int_0^T|\nabla\theta_i(X(v))|^2\,dv
 \le T\|w\|_\infty\cE_g(\theta_i)<\infty,
\]
because \(\|\widetilde S_vw\|_\infty\le\|w\|_\infty\).
This proves \eqref{entr:angular-sde}.  Countability permits a common
null set for all indices.
\end{proof}

We now cancel the common head projection in the angular martingale and
the common angular drift with a single stochastic integrating factor.
On the process of Lemma~\ref{entr:angular-realization}, define the
positive scalar \(\mathcal U(t)\) by
\begin{equation}\label{entr:scalar-rescaling}
 \mathcal U(0)=\frac{R(0)}{\sqrt{M(X(0))}},\qquad
 d\mathcal U
 =\mathcal U\left\{
     \frac{\sqrt2}{R}\sum_{j\le m}\Theta_j\,dW_j
       +\left(\frac{m-1}{R^2}-\overline r\right)dt\right\}.
\end{equation}
This is the stochastic exponential with the displayed drift and initial
value.  On every compact time interval \(R\) has a strictly positive
minimum, so all its coefficients are locally bounded along the path;
the exponential is finite and strictly positive.  Put
\begin{equation}\label{entr:YQ}
 Q=\frac{\mathcal U}{R},\qquad f=Q^2,\qquad
 Y_i=\mathcal U\Theta_i=QX_i,\qquad y_i=Y_i(0).
\end{equation}
In the product rule for \(\mathcal U\Theta_i\), the quadratic
covariation contributes
\(2\mathcal U\Theta_i(\ind_{\{i\le m\}}-1)R^{-2}dt\).
Adding this to the two drifts in
\eqref{entr:angular-sde}--\eqref{entr:scalar-rescaling} cancels all
terms containing \(R^{-2}\) or \(\overline r\).  Their martingale
terms cancel the head projection.  Thus
\begin{equation}\label{entr:common-sde}
 dY_i=-r_iY_i\,dt+\sqrt2Q\,dW_i,\qquad
 Y_i(t)=e^{-r_it}y_i+\sqrt2\int_0^t e^{-r_i(t-v)}Q(v)\,dW_i(v).
\end{equation}
Both \(Q\) and \(f\) are positive continuous processes, including at
elapsed time zero.  The normalization at zero gives
\begin{equation}\label{entr:seed-weights}
 p_i:=\frac{y_i^2}{2r_i}\ge0,\qquad
 \sum_i p_i=1,\qquad
 y_i^2=\frac{X_i(0)^2}{M(X(0))}.
\end{equation}
In particular \eqref{entr:escape} says that these weights leave every
finite set in probability when the starting time decreases along its
specified sequence.

\subsection{The scalar constraint equation}

The coordinates in \eqref{entr:common-sde} are driven by the same
factor \(Q\).  The quadratic constraint on \(X\) determines that
factor through a scalar Volterra equation.  To identify the equation
before taking an infinite limit, view \(f\) and the time kernels causally,
by extending them by zero to negative elapsed times, and fix \(N\).  Set
\[
 E_i(t)=e^{-2r_it}\quad(t\ge0),\qquad
 L_N(t)=\sum_{i\le N}r_i^{-1}E_i(t),\qquad
 R_N(x)=D-\sum_{i\le N}(x_i^2-r_i^{-1}).
\]
Itô's formula in \eqref{entr:common-sde} gives
\[
 Y_i(t)^2=E_i(t)y_i^2+2(E_i*f)(t)
       +2\sqrt2\int_0^t E_i(t-v)Q(v)Y_i(v)\,dW_i(v).
\]
Since \(\sum_{i\le N}Y_i^2=f\sum_{i\le N}X_i^2\) and
\[
 \sum_{i\le N}X_i^2
   =D+\sum_{i\le N}r_i^{-1}-R_N(X),
\]
rearrangement gives, as an identity of causal distributions,
\begin{equation}\label{entr:finite-constraint}
 Df+(L_N*f)'=
 \sum_{i\le N}E_i y_i^2
 +2\sqrt2\sum_{i\le N} E_i*(QY_i\,dW_i)+fR_N(X).
\end{equation}
Here the derivative includes the mass at elapsed time zero:
\begin{equation}\label{entr:initial-delta}
 L_N'=\left(\sum_{i\le N}r_i^{-1}\right)\delta_0
                       -2\sum_{i\le N}E_i .
\end{equation}
Thus the coefficient of \(f\) contributed by the delta mass is retained
in \eqref{entr:finite-constraint}.  This identity is exact for finite
\(N\).  The remaining passage must remove \(fR_N(X)\) and replace
\(L_N\) by
\[
 L(t)=\sum_i r_i^{-1}E_i(t),\qquad t>0.
\]
The corresponding limiting operator on \(f\) is convolution with
\(D\delta_0+L'\).  The next lemma constructs its positive inverse
kernel and the bounds needed for the stochastic terms.

\begin{lemma}[Resolvent kernels]\label{entr:kernels}
For \(L\) defined above, let
\[
 \widehat K(z)=\frac{1}{D+z\widehat L(z)}\quad(z>0),
\]
where a hat denotes the Laplace transform and every kernel is zero on
negative times.  The function \(L\) is locally integrable, and
\(\widehat K\) is the Laplace transform of a nonnegative, decreasing
function \(K\), smooth on \((0,\infty)\).  For \(0<t\le1\),
\begin{equation}\label{entr:K-bounds}
 c_g t^{-1/2}\le K(t)\le C_g t^{-1/2},
 \qquad |K'(t)|\le C_g t^{-3/2}.
\end{equation}
Set \(K_i=K*E_i\).  Then
\begin{align}
 K_i(t)&\le C_g\frac{\sqrt t}{1+r_it},&
 |K_i'(t)|&\le C_g\frac{t^{-1/2}}{1+r_it},             \label{entr:Ki-bounds}\\
 2r_iK_i(t)&\longrightarrow K(t)                       \label{entr:Ki-limit}
\end{align}
uniformly on each compact subinterval of \((0,1]\) as \(i\to\infty\).
For \(t>0\), \(h\ge0\), \(t+h\le1\), and \(0\le\xi<1/4\),
\begin{align}
 \sup_i r_iK_i(t)^2&\le C_g,&
 \sup_i r_i|K_i(t+h)-K_i(t)|^2
   &\le C_g\min\{1,h^2/t^2\},                          \label{entr:Ki-sup}\\
 \sum_i(1+r_i)^{2\xi}K_i(t)^2
   &\le C_{g,\xi}t^{-1/2-2\xi},                       \label{entr:Ki-sum}\\
 \sum_i|K_i(t+h)-K_i(t)|^2
   &\le C_g t^{-1/2}\min\{1,h^2/t^2\}.                \label{entr:Ki-increment}
\end{align}
As causal distributions,
\begin{equation}\label{entr:kernel-inverse}
 K*(D\delta_0+L')=\delta_0.
\end{equation}
\end{lemma}

\begin{proof}
The local integrability follows directly from
\[
 \int_0^T L(t)\,dt
 =\frac12\sum_i r_i^{-2}(1-e^{-2r_iT})<\infty.
\]
For \(L_N=\sum_{i\le N}r_i^{-1}E_i\), put
\[
 F_N(z)=D+z\widehat L_N(z)
       =D+\sum_{i\le N}\frac{z}{r_i(z+2r_i)}.
\]
The rational function \(F_N\) is strictly increasing between its
negative poles, since
\(F_N'(z)=\sum_{i\le N}2(z+2r_i)^{-2}>0\).
Its zeros interlace those poles and are strictly negative, using \(D>0\).
Its reciprocal is therefore a positive constant plus a sum of positive
multiples of \((z+a)^{-1}\), \(a>0\): the residue at a zero is
\(1/F_N'(-a)>0\).  Repeated rates can be grouped before this argument.

Write these sums as Stieltjes transforms of positive measures.  Their
measures weighted by \((1+a)^{-1}\), with the constant as mass at the
point at infinity, have total mass \(F_N(1)^{-1}\le D^{-1}\).
Weak compactness on \([0,\infty]\) and monotone convergence of
\(F_N(z)\) give
\[
 \frac1{D+z\widehat L(z)}
   =c+\int_{[0,\infty)}\frac{\nu(da)}{z+a},
 \qquad \int\frac{\nu(da)}{1+a}<\infty.
\]
By \eqref{entr:rates}, splitting the sum at \(r_i=z\) gives
\begin{equation}\label{entr:laplace-scale}
 D+z\widehat L(z)\asymp_g \sqrt z\qquad(z\ge1).
\end{equation}
For sufficiently large \(z\), the part with \(r_i\le z\) is
comparable to \(\sum_{r_i\le z}r_i^{-1}\asymp\sqrt z\), and the
other part to \(z\sum_{r_i>z}r_i^{-2}\asymp\sqrt z\).
Positivity and continuity adjust the constants on the remaining
compact interval \(z\ge1\).
Hence the transform tends to zero at infinity and \(c=0\).
It is bounded as \(z\downarrow0\), so \(\nu\) has no atom at zero.
The function
\[
 K(t)=\int_{(0,\infty)}e^{-at}\nu(da)
\]
has the required Laplace transform by Tonelli's theorem.  The weighted
integrability of \(\nu\) implies smoothness for \(t>0\), and the
representation makes \(K\) positive and decreasing.

The upper bound in \eqref{entr:K-bounds} follows from
\[
 e^{-1}tK(t)\le\int_0^t e^{-v/t}K(v)\,dv
 \le\widehat K(1/t)\le C_g\sqrt t.
\]
For the lower bound choose a fixed small \(\eta>0\).  When \(t\le\eta\),
\eqref{entr:laplace-scale} bounds
\(\widehat K(\eta/t)\) below by \(c_g\sqrt{t/\eta}\), whereas the
upper bound just proved bounds its integral over \([0,t]\) by
\(C_g\sqrt t\).  Choose \(\eta\) so that the first constant is larger
than twice the second.  The integral over \([t,\infty)\) is at most
\(\eta^{-1}t e^{-\eta}K(t)\), proving the lower bound for \(t\le\eta\).
Positivity and monotonicity adjust the constant on \([\eta,1]\).
Finally
\[
 |K'(t)|=\int a e^{-at}\nu(da)
 \le \frac{C}{t}K(t/2)\le C_g t^{-3/2}.
\]

For \(r_it\le1\), integration of \(K(v)\le C_gv^{-1/2}\)
gives the first bound in \eqref{entr:Ki-bounds}.  For \(r_it>1\),
write \(K_i(t)=\int_0^t K(t-v)e^{-2r_iv}\,dv\) and split at \(t/2\).
The first half is bounded by \(C_g/(r_i\sqrt t)\); the second by
\(C_g\sqrt t\,e^{-r_it}\), which has the same bound.  For the
derivative use \(K_i'=K-2r_iK_i\).  When \(r_it\le1\) the preceding
bounds suffice.  When \(r_it>1\), write
\[
 K_i'(t)=K(t)e^{-2r_it}
  -2r_i\int_0^t e^{-2r_iv}\bigl(K(t-v)-K(t)\bigr)\,dv.
\]
On \(v\le t/2\) the derivative bound on \(K\) makes the integrand
difference at most \(C_g v t^{-3/2}\); on \(v>t/2\) the exponential
factor pays the integral of \(K\).  This gives
\(|K_i'(t)|\le C_g/(r_it^{3/2})\).  It also proves
\eqref{entr:Ki-limit}, uniformly away from zero.

The remaining estimates follow by summing
\eqref{entr:Ki-bounds}.  For example \eqref{entr:rates} gives
\[
 \sum_i\frac{(1+r_i)^{2\xi}}{(1+r_it)^2}
 \le C_{g,\xi}t^{-3/2-2\xi}\quad(0\le\xi<1/4).
\]
The restriction on \(\xi\) is exactly integrability at infinity after
the change of variable \(v=r_it\): the comparison integral there is
\(\int_0^\infty v^{1/2+2\xi}(1+v)^{-2}\,dv\).
This proves \eqref{entr:Ki-sum}.  If \(h<t\), integrate the bound on
\(K_i'\) between \(t\) and \(t+h\), and then use the same sum with
\(\xi=0\); this gives \eqref{entr:Ki-increment} and
\eqref{entr:Ki-sup}.  If \(h\ge t\), bound the two values separately.
Finally the locally integrable function
\[
 D\int_0^tK(v)\,dv+(K*L)(t)
\]
has Laplace transform \(1/z\), and hence equals \(1\) for almost
every \(t>0\).  Its causal distributional derivative gives
\eqref{entr:kernel-inverse}.
\end{proof}

For the normalized seed weights \(p_i=y_i^2/(2r_i)\), define, for
\(0<t\le1\), the candidate seed response
\[
 A_y(t)=\sum_iK_i(t)y_i^2=\sum_i p_i\,2r_iK_i(t).
\]
The bound \(2r_iK_i(t)\le C_g t^{-1/2}\), together with
\(\sum_i p_i=1\), proves that this series converges uniformly on
each compact subinterval of \((0,1]\).  Its sum \(A_y\) is continuous
there and satisfies \(0\le A_y(t)\le C_g t^{-1/2}\).
The high-index limit \(2r_iK_i\to K\) is the ingredient used below
to prove \(A_y\to K\) in probability, uniformly on compact subintervals
of \((0,1]\), when the weights leave every fixed finite set.
The bound \(K(t)\asymp_g t^{-1/2}\) determines the scale used in
the later stopped comparison of \(f=Q^2\).  The next lemma defines
the infinite stochastic convolution and controls it uniformly over
normalized seeds under the corresponding upper bound
\(|q(t)|\le C_0t^{-1/4}\).

\begin{lemma}[Stochastic kernel bounds]\label{entr:stochastic-kernels}
Let \((W_i)\) be independent Brownian motions in a filtration, let
\((y_i)\) be measurable at time zero with
\(\sum_i y_i^2/(2r_i)=1\), and let \(q\) be progressively measurable
with \(|q(v)|\le C_0v^{-1/4}\) for \(0<v\le1\).  Define
\begin{align}
 Z_i^q(t)&=\sqrt2\int_0^t e^{-r_i(t-v)}q(v)\,dW_i(v),\notag\\
 \mathsf G_q^y(t)&=2\sqrt2\sum_i\int_0^t
                K_i(t-v)q(v)e^{-r_iv}y_i\,dW_i(v),\label{entr:stochastic-parts}\\
 \mathsf N_q(t)&=2\sqrt2\sum_i\int_0^t
                K_i(t-v)q(v)Z_i^q(v)\,dW_i(v).\notag
\end{align}
The two series converge in probability in \(C([0,1])\), with value
zero at zero.  For every finite \(j\ge2\), their partial sums and their
limits satisfy
\begin{align}
 \|Z_i^q(t)\|_{L^j}^2
     &\le C_{g,C_0,j}\frac{\sqrt t}{1+r_it},           \label{entr:Z-moment}\\
 \|\mathsf G_q^y(t)\|_{L^j}+\|\mathsf N_q(t)\|_{L^j}
     &\le C_{g,C_0,j}t^{1/4},                         \label{entr:stoch-moment}\\
 \|\mathsf G_q^y(t+h)-\mathsf G_q^y(t)\|_{L^j}
  +\|\mathsf N_q(t+h)-\mathsf N_q(t)\|_{L^j}
     &\le C_{g,C_0,j}h^{1/4}\quad(0\le t<t+h\le1).
                                                               \label{entr:stoch-increment}
\end{align}
The supremum on \([0,1]\) of either limit has every finite moment,
bounded by a constant depending only on \(g,C_0\), and that order.
\end{lemma}

\begin{proof}
The bracket estimate and the Burkholder--Davis--Gundy inequality give
\[
 \|Z_i^q(t)\|_{L^j}^2
 \le C_{C_0,j}\int_0^t e^{-2r_i(t-v)}v^{-1/2}\,dv
 \le C_{C_0,j}\frac{\sqrt t}{1+r_it}.
\]
For the last inequality split at \(t/2\); on the second half use
\(v^{-1/2}\le C t^{-1/2}\), and on the first use
\(e^{-2r_i(t-v)}\le e^{-r_it}\).

Write \(p_i=y_i^2/(2r_i)\).  The bracket of the first series, at a
fixed time \(t\), is bounded pathwise by
\[
 C_{C_0}\int_0^t v^{-1/2}
       \sum_i p_i r_iK_i(t-v)^2\,dv
 \le C_{g,C_0}\sqrt t,
\]
by \eqref{entr:Ki-sup} and \(\sum_i p_i=1\).
This proves its bound in \eqref{entr:stoch-moment}.  In its increment,
the old portion of the bracket is at most
\[
 C_{g,C_0}\int_0^t v^{-1/2}
              \min\{1,h^2/(t-v)^2\}\,dv
 \le C_{g,C_0}\sqrt h ,
\]
and the new portion is bounded by
\(C\int_t^{t+h}v^{-1/2}\,dv\le C\sqrt h\).
For the displayed integral, if \(t\le2h\) discard the minimum; if
\(t>2h\), split at \(t-h\) and at \(t/2\).  These bounds prove the
first increment estimate.

For the second series, the martingale moment inequality followed by
Minkowski's inequality bounds its \(L^j\) norm squared by
\[
 C_j\int_0^t\sum_i K_i(t-v)^2
                      \|q(v)Z_i^q(v)\|_{L^j}^2\,dv
 \le C_{g,C_0,j}\int_0^t (t-v)^{-1/2}\,dv
 \le C_{g,C_0,j}\sqrt t.
\]
We used \eqref{entr:Z-moment}, the deterministic bound on \(q\), and
\eqref{entr:Ki-sum} with \(\xi=0\).
For the increment use \eqref{entr:Ki-increment} on the old portion:
its squared norm is bounded by
\[
 C_{g,C_0,j}\int_0^t (t-v)^{-1/2}
                \min\{1,h^2/(t-v)^2\}\,dv
 \le C_{g,C_0,j}\sqrt h .
\]
The new portion has the same bound by integrating
\((t+h-v)^{-1/2}\) from \(t\) to \(t+h\).

All these estimates hold for partial sums.  At each fixed time their
tails tend to zero in \(L^j\): for the first series use
\(\sum_{i>N}p_i r_iK_i(t-v)^2\to0\) and its constant bound inside
the bracket; for the second use the summable bound just displayed.
Dominated convergence applies in both cases.  Taking \(j>4\),
\eqref{entr:stoch-increment} gives uniformly small moduli of
continuity, by a dyadic chaining estimate (equivalently the
Kolmogorov continuity criterion).  Fixed-time convergence on a finite
grid then proves convergence in probability in \(C([0,1])\).
The same chaining has a convergent series
\(\sum_{k\ge0}2^{k/j}2^{-k/4}\), and bounds the \(L^j\) norm of the
supremum for \(j>4\).  Taking a higher order gives the asserted bound
for every finite order.
\end{proof}

\begin{lemma}[The renormalized Volterra equation]\label{entr:volterra}
Let the processes be as in \eqref{entr:YQ}--\eqref{entr:common-sde},
and let \(A_y\) be the candidate seed response defined above.
Almost surely, for \(0<t\le1\),
\begin{equation}\label{entr:volterra-equation}
 f(t)=A_y(t)+2\sqrt2\sum_i\int_0^t
               K_i(t-v)Q(v)Y_i(v)\,dW_i(v).
\end{equation}
The stochastic series is defined by localization and has a continuous
version on every compact subinterval of \((0,1]\).
\end{lemma}

\begin{proof}
We now justify the infinite passage in \eqref{entr:finite-constraint}.
The constraint says \(R_N(x)=\sum_{i>N}(x_i^2-r_i^{-1})\)
\(\Pi\)-almost surely.  The conditioned edge law describes the tail,
after any sufficiently large fixed head, by a bounded density relative
to its product Gaussian law
\cite[Proposition~2.3]{FUA}.
Consequently, for \(N\) beyond that head,
\begin{equation}\label{entr:tail-residual}
 \Pi[|R_N|]\le C_g\left(\sum_{i>N}r_i^{-2}\right)^{1/2}
       \longrightarrow0.
\end{equation}
Indeed the Gaussian tail is centered and has variance
\(2\sum_{i>N}r_i^{-2}\).

The density of \(X(t)\) is bounded by \(\|w\|_\infty\) for all
elapsed \(t\ge0\).  The constraint therefore holds for almost every
elapsed time almost surely, and \eqref{entr:tail-residual} gives
\[
 \E\int_0^T|R_N(X(v))|\,dv
       \le T\|w\|_\infty\Pi[|R_N|]\longrightarrow0 .
\]
The continuous process \(f\) is bounded on \([0,T]\) almost surely.
First restrict to a bound for that supremum, and then let the bound
increase.  It follows that
\(fR_N(X)\to0\) in probability in \(L^1(0,T)\).
Also
\begin{equation}\label{entr:L-tail}
 \|L-L_N\|_{L^1(0,T)}
 \le\frac12\sum_{i>N}r_i^{-2}\longrightarrow0 .
\end{equation}
Convolving \eqref{entr:finite-constraint} with \(K\), its left side
therefore tends to \(f\) in local distributions by
\eqref{entr:kernel-inverse}.  One can see this directly by writing
the difference as the derivative of
\(K*(L_N-L)*f\); its undifferentiated factor tends to zero in
\(L^1_{\mathrm{loc}}\) by \eqref{entr:L-tail}.  The convolved residual
also tends to zero in probability in \(L^1_{\mathrm{loc}}\), since
\(K\) is locally integrable.

The convolved deterministic terms are
\(\sum_{i\le N}K_i(t)y_i^2\), which converge uniformly on compact
subintervals of \((0,1]\) to \(A_y(t)\) by its definition and bound
above.

For a finite stochastic sum, stochastic Fubini gives
\[
 K*\bigl(E_i*(QY_i\,dW_i)\bigr)(t)
       =\int_0^t K_i(t-v)Q(v)Y_i(v)\,dW_i(v).
\]
It can first be used with \(QY_i\) stopped at a finite bound.
For such a bound a sufficient stochastic Fubini integral is at most
\[
 C\int_0^t K(t-s)
       \left(\int_0^s E_i(s-v)^2\,dv\right)^{1/2}ds
 \le C\int_0^t K(t-s)\sqrt s\,ds<\infty.
\]
Removing that bound by localization gives the displayed identity as
local distributions.

To specify the limit without a formal exchange of infinite stochastic
sums, for \(C_0>0\) set
\[
 q_{C_0}(v)=Q(v)\wedge(C_0v^{-1/4}),\qquad
 \zeta_{C_0}=\inf\{v>0:Q(v)>C_0v^{-1/4}\}.
\]
By Lemma~\ref{entr:stochastic-kernels}, the partial sums of
\(\mathsf G_{q_{C_0}}^y+\mathsf N_{q_{C_0}}\) converge in probability
in \(C([0,1])\) on the whole probability space.  Before
\(\zeta_{C_0}\), \eqref{entr:common-sde} gives
\(Y_i(v)=e^{-r_iv}y_i+Z_i^{q_{C_0}}(v)\), so these clipped partial
sums agree there with the convolved finite stochastic sums.

Fix a deterministic \(C_0>0\) and a rational \(a\in(0,1]\).
On the whole probability space, the clipped partial sums converge in
probability in \(C([0,a])\), and the convolved residual
\(K*(fR_N(X))\) converges in probability to zero in \(L^1(0,a)\)
as shown above.  We may therefore choose a deterministic subsequence
\(N_j\uparrow\infty\) along which both convergences hold almost surely.
Only then restrict to the event \(\{a<\zeta_{C_0}\}\).  On this event,
passage in \eqref{entr:finite-constraint} proves
\eqref{entr:volterra-equation} as a distribution on \((0,a)\).
Take \(C_0\) through the positive integers and \(a\) through the
rationals in \((0,1]\), and intersect the corresponding whole-space
probability-one convergence events.  The continuous versions then give
equality for \(0<t<\min\{\zeta_{C_0},1\}\), and by continuity at a
positive hitting time \(\zeta_{C_0}\le1\) and at \(t=1\) when
\(\zeta_{C_0}>1\).
The positive continuous \(Q\) is finite at zero, and
\(\sup_{v\le1}v^{1/4}Q(v)<\infty\) almost surely.  A sufficiently
large integer \(C_0\) therefore covers the whole interval.  This proves the
claimed identity and its precise interpretation.
\end{proof}

We next isolate the part of the equation that remembers the initial
state.  It becomes independent of that state when the weights
\eqref{entr:seed-weights} escape.

\begin{lemma}[Escaping seeds]\label{entr:escaping-seeds}
Suppose that, along a sequence of initial states possibly random,
\(\sum_i p_i=1\) and \(p_i\to0\) in probability for every fixed \(i\).
Then
\begin{equation}\label{entr:A-limit}
 A_y\longrightarrow K
 \quad\text{in probability, uniformly on compact subintervals of \((0,1]\)}.
\end{equation}
For each fixed \(t>0\), the seed stochastic term
\(\mathsf G_q^y(t)\) tends to zero in \(L^2\), uniformly over all
adapted \(q\) satisfying \(|q(v)|\le C_0v^{-1/4}\).
\end{lemma}

\begin{proof}
By \eqref{entr:Ki-limit},
\(2r_iK_i\to K\) uniformly on a fixed compact positive interval,
and these functions are uniformly bounded there.  In the convex
combination \(A_y=\sum_i p_i\,2r_iK_i\), the contribution of a
fixed finite set tends to zero in probability and the contribution
of its complement is uniformly close to \(K\).  This proves
\eqref{entr:A-limit}.

For the stochastic assertion, its squared norm is at most
\[
 C_{C_0}\int_0^t v^{-1/2}\,
       \E\!\left[\sum_i p_i r_i K_i(t-v)^2\right]\,dv .
\]
For every \(v<t\), \eqref{entr:Ki-bounds} implies
\(r_iK_i(t-v)^2\to0\), whereas \eqref{entr:Ki-sup} bounds it
uniformly.  Escape of \(p_i\) makes the expression in brackets tend
to zero in probability, and its uniform bound gives convergence of
its expectation.  Dominated convergence against \(v^{-1/2}\)
proves the assertion.
\end{proof}

\subsection{Coupling entrances}

We now prove the deterministic uniqueness and boundary assertions of
Theorem~\ref{thm:entrance-selection}.  Work in unit time
\eqref{entr:unit-time}.  Let two proposed families be given, with
starting sequences \(u_{\ell,k}\downarrow0\), \(\ell=1,2\).
Start the diffusion of Lemma~\ref{entr:angular-realization} from
\(\widetilde h^{\,\ell}_{u_{\ell,k}}\Pi\), and construct
\((Q_\ell,f_\ell,Y_{\ell,i},y_{\ell,i})\) as above.  Their seed weights
escape in probability by \eqref{entr:escape} and
\eqref{entr:seed-weights}.  The constants below depend on \(g\), and
are uniform in \(k\) and in the two families.

We may use the same Brownian motions for the two diffusions in elapsed
time even though strong existence has not been asserted.  Here is the
coupling argument.  Each joint law of the Brownian path \(W\) and all
its other continuous path variables is a probability on a standard
Borel path space.  Disintegrate it over \(W\), and, given one Brownian
path, sample the two other path variables independently from these
conditional laws.  Brownian motion remains Brownian in the joint
filtration.  To see this, if \(\Phi_\ell\) is a bounded observable of the
past of process \(\ell\) through time \(t\), independence of future
Brownian increments in that process implies
\[
 \E[\Phi_\ell\mid W\text{ on }[0,1]]
   =\E[\Phi_\ell\mid W\text{ on }[0,t]]
 \quad\text{almost surely}.
\]
Test this equality first against products of functions of the Brownian
past and future increments, and then use a monotone class argument.
For a product \(\Phi_1\Phi_2\), conditional independence in the coupling
makes its conditional expectation the product of two functions of the
Brownian past.  A second monotone class argument proves independence
of future increments from the joint past.  Completion and
right-continuous augmentation preserve this property.  Finally,
stochastic integrals are limits in probability of adapted simple
sums.  Their laws in each marginal are unchanged by the coupling,
and they are the same integrals in this joint Brownian filtration.
Thus all the preceding stochastic identities hold on the coupling.

\paragraph{Stops with uniform failure bounds.}
Choose \(C_0\) so large that
\[
 C_0^2>1+\sup_{0<t\le1}\sup_y\sqrt t\,A_y(t),
\]
where the supremum is over the normalized seeds in
\eqref{entr:seed-weights}; it is finite by the bound on \(A_y\)
following Lemma~\ref{entr:kernels}.  For each solution use the clipped process
\(q_\ell^{\,0}(v)=Q_\ell(v)\wedge(C_0v^{-1/4})\).
Let \(k_0=\inf_{0<t\le1/2}\sqrt t\,K(t)>0\), and choose
\(c\in(0,k_0/4)\).  Lemma~\ref{entr:escaping-seeds} permits a
deterministic sequence \(\eta_k\downarrow0\), as slow as necessary,
such that with probability tending to one both solutions satisfy
\begin{equation}\label{entr:seed-lower}
 \sqrt t\,A_{y_\ell}(t)\ge2c
             \qquad(\eta_k\le t\le1/2).
\end{equation}
This follows by a diagonal choice from uniform convergence on each
fixed compact interval.

Stop each solution when \(\sqrt t\,f_\ell(t)\) first reaches the
upper threshold \(C_0^2\); it starts below that threshold since
\(f_\ell\) is finite at zero.  Starting at \(\eta_k\), also stop
when \(\sqrt t\,f_\ell(t)\le c\).  Before the upper stop,
Lemma~\ref{entr:volterra} reads
\[
 f_\ell=A_{y_\ell}
       +\mathsf G_{q_\ell^{\,0}}^{y_\ell}
       +\mathsf N_{q_\ell^{\,0}}.
\]
It holds also at a positive first hit by continuity.  On the event
\eqref{entr:seed-lower}, either threshold hit by time \(\delta\le1/2\)
requires the supremum of the stochastic term on \([0,1]\) to be
at least a fixed multiple of \(\delta^{-1/2}\).
Lemma~\ref{entr:stochastic-kernels} and Markov's inequality therefore
bound its probability by \(C_j\delta^{j/2}\) for any fixed
sufficiently large \(j\), plus a term tending to zero with \(k\).

Fix \(0<\xi<1/4\).  Add a stop for each solution at the infimum,
over \(i\ge1\), of the first times
\[
 |Z_i^{q_\ell^{\,0}}(t)|\ge(1+r_i)^\xi.
\]
This costs at most a positive power of \(\delta\) by time \(\delta\).
Indeed if \(N_i(t)=\sqrt2\int_0^t q_\ell^{\,0}(v)\,dW_i(v)\),
integration by parts gives
\[
 Z_i^{q_\ell^{\,0}}(t)
 =N_i(t)-r_i\int_0^t e^{-r_i(t-v)}N_i(v)\,dv,
 \qquad
 \sup_{t\le\delta}|Z_i^{q_\ell^{\,0}}(t)|
       \le2\sup_{t\le\delta}|N_i(t)|.
\]
The last supremum has \(L^j\) norm at most \(C_j\delta^{1/4}\),
uniformly in \(i\).  Choose \(j>3/(2\xi)\), so
\(\sum_i(1+r_i)^{-j\xi}<\infty\) by \eqref{entr:rates}, and sum
Markov's inequality over \(i\).  The resulting bound is
\(C_j\delta^{j/4}\).  A countable infimum of these stopping times is
a stopping time in the right-continuous filtration.  If an infimum is
not attained, an individual stop occurs by \(2\delta\) whenever the
infimum is at most \(\delta>0\); this only changes the constant.

Let \(\tau_k\) be the minimum of all these stops for both solutions,
capped at \(1/2\).  There are constants \(b>0,C<\infty\), independent
of \(k\), such that, for sufficiently small fixed \(\delta>0\),
\begin{equation}\label{entr:stop-probability}
 \limsup_{k\to\infty}\Prob(\tau_k\le\delta)\le C\delta^b.
\end{equation}
Here and below the probability is over the initial draws and the
coupled diffusions, with \(g\) fixed.

\paragraph{Contraction of the common factors.}
Use the common stopped integrands
\[
 q_\ell(v)=Q_\ell(v)\ind_{\{v<\tau_k\}},\qquad
 d_k(v)=\sqrt v\,\E|q_1(v)-q_2(v)|^2.
\]
They satisfy \(|q_\ell(v)|\le C_0v^{-1/4}\), so \(0\le d_k(v)\le C\).
Until \(\tau_k\), the stopped versions of
\eqref{entr:volterra-equation} agree with \(f_\ell\), and
\(|Z_i^{q_\ell}(v)|\le(1+r_i)^\xi\).
For \(v\ge\eta_k\), the lower threshold gives
\begin{equation}\label{entr:sqrt-difference}
 d_k(v)\le C v\,\E\left[
       \ind_{\{v<\tau_k\}}|f_1(v)-f_2(v)|^2\right].
\end{equation}
This is the identity
\((\sqrt{f_1}-\sqrt{f_2})^2=(f_1-f_2)^2/
(\sqrt{f_1}+\sqrt{f_2})^2\), using
\(f_\ell(v)\ge c v^{-1/2}\) before the stop.

Let
\[
 F_\ell(v)=A_{y_\ell}(v)
       +\mathsf G_{q_\ell}^{y_\ell}(v)+\mathsf N_{q_\ell}(v).
\]
It agrees with \(f_\ell\) for \(v<\tau_k\).  At a fixed \(v>0\),
Lemma~\ref{entr:escaping-seeds} and its uniform bounds give
\begin{equation}\label{entr:seed-vanishes}
 \E|A_{y_1}(v)-A_{y_2}(v)|^2
  +\E|\mathsf G_{q_1}^{y_1}(v)|^2
  +\E|\mathsf G_{q_2}^{y_2}(v)|^2\longrightarrow0 .
\end{equation}
For the deterministic term, convergence in probability becomes
\(L^2\) convergence because \(|A_y(v)|\le C_gv^{-1/2}\).

It remains to compare \(\mathsf N_{q_1}\) and \(\mathsf N_{q_2}\).
In its integrand write
\[
 q_1Z_i^{q_1}-q_2Z_i^{q_2}
   =(q_1-q_2)Z_i^{q_1}+q_2(Z_i^{q_1}-Z_i^{q_2}).
\]
The first term is nonzero only before the common stop, where
\(|Z_i^{q_1}|\le(1+r_i)^\xi\).  For the second, the common
Brownian motions and Itô isometry give
\[
 \E|Z_i^{q_1}(w)-Z_i^{q_2}(w)|^2
 =2\int_0^w e^{-2r_i(w-a)}a^{-1/2}d_k(a)\,da .
\]
Using \eqref{entr:Ki-sum} with \(\xi\) for the first term and
with zero for the second yields
\begin{align}
 \E|\mathsf N_{q_1}(v)-\mathsf N_{q_2}(v)|^2
 &\le C\int_0^v (v-w)^{-1/2-2\xi}w^{-1/2}d_k(w)\,dw\notag\\
 &\quad+C\int_0^v (v-w)^{-1/2}w^{-1/2}
                 \int_0^w a^{-1/2}d_k(a)\,da\,dw .
                                                        \label{entr:contraction-integrals}
\end{align}
All sums and integrals here are justified by Tonelli's theorem and
the displayed finite upper bounds.  In the second line we used the
deterministic bound \(q_2(w)^2\le C_0^2w^{-1/2}\) and discarded
the decaying exponential after Itô isometry.  No independence of
these integrands is required.

For \(v>0\) define \(d_\infty(v)=\limsup_k d_k(v)\).
For each fixed \(v\) we eventually have \(v\ge\eta_k\).
Combine \eqref{entr:sqrt-difference}--\eqref{entr:contraction-integrals}.
The bounds \(d_k\le C\) allow the upper limit to pass under the
integrals by dominated convergence applied to upper envelopes.  The
kernels are integrable because \(\xi<1/4\).  If
\(D_\delta=\sup_{0<v\le\delta}d_\infty(v)\), the two beta integrals
then give
\[
 d_\infty(v)
 \le C\bigl(v^{1-2\xi}+v^{3/2}\bigr)D_\delta
 \qquad(0<v\le\delta).
\]
For the first power use
\(\int_0^v(v-w)^{-1/2-2\xi}w^{-1/2}\,dw=C_\xi v^{-2\xi}\);
for the second use
\(\int_0^w a^{-1/2}\,da=2\sqrt w\).
Choose \(\delta>0\) small enough that
\(C(\delta^{1-2\xi}+\delta^{3/2})<1\).  Taking the supremum proves
\begin{equation}\label{entr:q-convergence}
 d_k(v)\longrightarrow0\qquad(0<v\le\delta).
\end{equation}

At a fixed such \(\delta\), dominated convergence in
\eqref{entr:contraction-integrals} and
\eqref{entr:seed-vanishes} now shows
\[
 \E\!\left[\ind_{\{\delta<\tau_k\}}
                   |f_1(\delta)-f_2(\delta)|^2\right]\longrightarrow0.
\]
For each fixed \(i\), the same conclusion holds for
\(|Y_{1,i}(\delta)-Y_{2,i}(\delta)|^2\).  Indeed
\(y_{\ell,i}\to0\) in probability by escape and
\(|y_{\ell,i}|^2\le2r_i\), while the squared norm of the difference
of the two \(Z_i\) terms is
\(2\int_0^\delta e^{-2r_i(\delta-v)}v^{-1/2}d_k(v)\,dv\to0\)
by \eqref{entr:q-convergence}.  On \(\{\delta<\tau_k\}\) the
denominators satisfy \(f_\ell(\delta)\ge c\delta^{-1/2}\), and
\[
 |Y_{\ell,i}(\delta)|
 \le\sqrt{2r_i}+(1+r_i)^\xi .
\]
Since \(X_{\ell,i}=Y_{\ell,i}/\sqrt{f_\ell}\), it follows that for
every finite \(j\) and every \(\varepsilon>0\),
\begin{equation}\label{entr:coordinate-coupling}
 \limsup_k\Prob\left(\max_{i\le j}
       |X_{1,i}(\delta)-X_{2,i}(\delta)|>\varepsilon\right)
       \le C\delta^b.
\end{equation}

\paragraph{Equality of the families.}
The law at elapsed time \(\delta\) of process \(\ell\) is
\(\widetilde h^{\,\ell}_{u_{\ell,k}+\delta}\Pi\).  It converges
in total variation to \(\widetilde h^{\,\ell}_\delta\Pi\):
the semigroup relation writes its density as
\(\widetilde S_{u_{\ell,k}}\widetilde h^{\,\ell}_\delta\), and
strong continuity in \(L^2(\Pi)\) implies \(L^1\) convergence.
The couplings of their two elapsed-time marginals are tight on
\(\R^\N\times\R^\N\), so take a weakly convergent subsequence.
Its limit is a coupling of the two prescribed marginals at
\(\delta\).  For the open event in
\eqref{entr:coordinate-coupling}, the portmanteau theorem gives the
same upper bound for the limiting coupling.  Let
\(\varepsilon\downarrow0\) and then \(j\uparrow\infty\).
The limiting coupling assigns probability at least \(1-C\delta^b\)
to equality of all coordinates.  Hence
\[
 \|\widetilde h^{\,1}_\delta\Pi
          -\widetilde h^{\,2}_\delta\Pi\|_{\TV}\le C\delta^b.
\]
The Markov semigroup contracts total variation.  For any fixed
\(v>\delta\) the same bound holds at \(v\), and sending
\(\delta\downarrow0\) proves equality of the two marginals at \(v\).
They are equal at every positive time, and hence their density classes
are equal.  This proves uniqueness.

\paragraph{The boundary at zero.}
Use one proposed family and the preceding stops.  On
\(\{\delta<\tau_k\}\),
\[
 |X_i(\delta)|
 =\frac{|Y_i(\delta)|}{\sqrt{f(\delta)}}
 \le \frac{\sqrt{2r_i}+(1+r_i)^\xi}{\sqrt c}\,\delta^{1/4}.
\]
Let the starting index \(k\to\infty\).  The elapsed-time law
converges in total variation to the prescribed marginal at \(\delta\),
and \eqref{entr:stop-probability} bounds the excluded probability by
\(C\delta^b\).  Thus each coordinate tends to zero in probability
under \(\widetilde h_\delta\Pi\) as \(\delta\downarrow0\).
For example the bounded product metric
\(\sum_i2^{-i}\min\{1,|x_i|\}\) tends to zero in probability:
truncate its sum and use the finitely many coordinate conclusions.
This proves \eqref{entr:zero-boundary} in unit time, and then in
the theorem's critical time by \eqref{entr:unit-time}.

\subsection{Measurability and Gaussian identification}

Return now to the theorem's critical time \(s=u/c_*\), so that the
evolution is by \(S_t^g=e^{-t c_*H_g}\).  It remains to obtain the family from \(g\)
itself and to check existence.
We give the measure-theoretic passage, since the escape estimate in
Proposition~\ref{prop:uniform-seed-escape} initially averages over the
represented environment as well as \(x\).

On any represented Gaussian limiting space write
\(\mu_s=h_s\Pi_g\).  At rational \(s>0\) these probability measures
are measurably encoded by their integrals against a countable
convergence-determining collection of bounded continuous cylinders.
Equivalently
\[
 \boldsymbol\mu=(\mu_s)_{s\in\mathbb Q_{>0}}
      \in\mathcal P(\R^\N)^{\mathbb Q_{>0}}
\]
is a random element of a standard Borel space.  The asserted
measurability follows from that encoding: a countable separating
collection embeds \(\mathcal P(\R^\N)\) injectively and measurably
in a countable product, and its inverse on the image is measurable.
By the convention in Proposition~\ref{prop:uniform-seed-escape}, the
escape quotient is jointly Borel on \(\R^\N\times\R^\N\).
For each fixed \(i\) and \(\varepsilon>0\), its threshold event is a
Borel subset of this parameter/state space, so the evaluation of that
event under \(\mu_s\) is a measurable function of \((g,\mu_s)\).

The annealed escape estimate permits a deterministic rational sequence
\(s_k\downarrow0\) for which, almost surely in the represented data,
\eqref{entr:escape} holds for every \(i,\varepsilon\).
To see this choose \(s_k<\min\{2^{-k},s_{k-1}/2\}\) rational,
successively, so that
\[
 \E\sum_{i\le k}
  \mu_{s_k}\{x_i^2/M_g(x)>1/k\}\le2^{-k}.
\]
Proposition~\ref{prop:uniform-seed-escape} allows this choice.
Tonelli's theorem shows that the sum of the displayed random
quantities over \(k\) is finite almost surely.  For fixed
\(i,\varepsilon>0\), eventually \(i\le k\) and \(1/k<\varepsilon\);
the corresponding probability in \eqref{entr:escape} consequently
tends to zero.  The positive-time semigroup and boundedness properties
hold on the same full-measure set by
Proposition~\ref{prop:gaussian-subsequence}.

The finite-coordinate formula \eqref{entr:cylinder-density} makes
\(g\mapsto\Pi_g\) a measurable probability kernel.  The edge
semigroup is also measurable in \(g\) in its countable varying-\(L^2\)
encoding.  To see the precise claim needed here, take a countable
smooth cylinder core.  Its Gram entries \(\Pi_g[\phi_i\phi_j]\) and
form entries \(\Pi_g[\nabla\phi_i\cdot\nabla\phi_j]\) are measurable.
For \(\lambda>0\) the variational formula
\[
 \langle\phi,(\lambda+H_g)^{-1}\phi\rangle_{\Pi_g}
 =\sup_v\{2\Pi_g[\phi v]-\lambda\Pi_g[v^2]-\cE_g(v)\},
\]
where \(v\) ranges over the rational span of that core, proves
measurability of the resolvent pairings; polarization gives mixed
pairings.  Choose the countable core sequence linearly independent
as cylinder functions, with dense real span.  The positive
finite-coordinate densities make every finite Gram matrix positive
definite.  The coefficients of the orthogonal projection of an encoded
\(L^2\) vector on a finite core span are its cylinder pairings
multiplied by the inverse Gram matrix; they are measurable.
These projections converge in \(L^2\), proving the corresponding
measurability for arbitrary encoded vectors.  Continuous functional
calculus then gives measurability of \(S_t^g\) in the same encoding,
jointly for \(t>0\); one can approximate \(e^{-tc_*H_g}\) by continuous
functions of \((1+H_g)^{-1}\), and use strong continuity in \(t\).
This is also the countable-core construction in the functional
identification proof for the stationary edge law
\cite[proof of Proposition~9.2]{FUA}.

The normalization, boundedness, evolution, and escape conditions are
properties of the countable code together with \(g\).
Normalization and bounded domination by \(\Pi_g\) can be checked on
a countable cylinder algebra, with an integer bound on each rational
compact time interval.  Evolution can be checked against the same
tests at rational times; it then extends to real times by strong
\(L^2\) continuity.  The preceding core construction makes the semigroup
pairings measurable in \(g\).  Thus these conditions, together with
the displayed countable escape tests, define a measurable full set
of codes on which the deterministic theorem applies.

Disintegrate the law of \((g,\boldsymbol\mu)\) over \(g\).
For almost every \(g\), two independent conditional samples of
\(\boldsymbol\mu\) both extend to the positive-time families just
considered and both pass \eqref{entr:escape}.  Fixed-\(g\) uniqueness
therefore makes these samples equal.  A probability measure on a
standard Borel space whose two independent draws agree almost surely
is a point mass (apply the assertion to a countable separating
collection of Borel sets).  Thus the conditional law of
\(\boldsymbol\mu\) given \(g\) is a point mass almost surely.
Its point depends measurably on \(g\): the map
\(z\mapsto\delta_z\) is a measurable injection into the space of
probability measures and has a measurable inverse on its image.
This constructs measurable \(\mu_s^g\) for rational \(s\), and
establishes existence.

We spell out the extension to density classes and real times.  A measurable
Radon--Nikodym theorem applied to the kernels
\(\mu_q^g\ll\Pi_g\) gives measurable representatives \(h_q^g(x)\)
for rational \(q\).

For real \(s>0\), choose any rational \(q<s\) and set
\begin{equation}\label{entr:real-extension}
 h_s^g=S_{s-q}^g h_q^g .
\end{equation}
The rational evolution relation makes this independent of \(q\).
Choosing \(q=q(s)\) by a fixed measurable rational rule and using
the semigroup measurability established above shows that the density
classes in \eqref{entr:real-extension} are measurable in \((g,s)\).
The corresponding measures form a measurable kernel; the measurable
Radon--Nikodym theorem once more gives jointly measurable
representatives \(h_s^g(x)\).  They are the \(L^2\)-continuous
extension furnished by the semigroup, with the stated normalization
and local boundedness.

For a second extracted Gaussian law, condition its code over the same
GOE law of \(g\).  Its conditional samples pass the criterion on some
sequence by the same argument, so fixed-\(g\) uniqueness makes its
point mass the one already constructed, outside a common null set.
Thus every subsequential Gaussian limit is determined by \(g\).
The locally uniform correlation convergence and relaxation in
Proposition~\ref{prop:gaussian-subsequence} now hold for this
\(h^g\), and the subsequence criterion proves the asserted joint
Gaussian convergence.  Finally, undoing \eqref{entr:unit-time}
returns the evolution \(h_{s+t}^g=e^{-t c_*H_g}h_s^g\).
The only time substitution in this section was the temporary variable
\(u=c_*s\); the coefficient in \(S_t^g\) remains \(c_*\).
\qed
